%% ---------------------------------------------------------------------------
%% Brain-Inspired Modular AI -- IEEE conference paper
%%
%% Build:
%%     python -m benchmark.run          produces results/results.json
%%     python -m benchmark.tables       produces paper/generated/*.tex
%%     pdflatex main && bibtex main && pdflatex main && pdflatex main
%%
%% NOTHING IN THE RESULTS SECTION IS TYPED BY HAND. Every number arrives via
%% \input{generated/...}, so the text cannot drift from results.json after a
%% rerun. Sentences that INTERPRET the numbers are marked \todo{} and must be
%% written only after looking at them -- see the note above Section VI.
%% ---------------------------------------------------------------------------

\documentclass[conference]{IEEEtran}
\IEEEoverridecommandlockouts

\usepackage{cite}
\usepackage{amsmath,amssymb}
\usepackage{graphicx}
\usepackage{booktabs}
\usepackage{xcolor}
\usepackage{url}

% Visible reminders that a claim still needs a number behind it.
\newcommand{\todo}[1]{\textcolor{red}{[TODO: #1]}}

\begin{document}

\title{Functionally Specialised Cortical Modules for Multimodal\\
       and Embodied Navigation}

\author{\IEEEauthorblockN{Rishiraj Chakraborty}
\IEEEauthorblockA{\textit{\todo{affiliation}} \\
\todo{city, country} \\
chakrabortyrishiraj09@gmail.com}}

\maketitle

% ---------------------------------------------------------------------------
\begin{abstract}
Conventional neural systems process heterogeneous inputs with a single
monolithic network. The mammalian cortex does not: distinct regions carry
distinct computational structure, a thalamic switchboard relays each sensory
modality to its own target, and damage to one region degrades some abilities
while sparing others. We ask whether that organisation is worth reproducing in
an artificial system, or whether it is only a biological accident.
%
We present a modular architecture in which four cortical modules are separate
networks with \emph{different inductive biases} -- locally connected
(occipital), relational self-attention (parietal), recurrent (temporal), and
convergent fusion (frontal) -- addressed by a modality-keyed thalamic relay and
gated by a basal-ganglia go/no-go mechanism. We evaluate on an embodied
navigation task in which vision, proprioceptive rangefinding and natural-language
instruction are each individually insufficient, plus two single-modality
control tasks.
%
Against parameter-matched baselines that remove exactly one property at a time
-- routing, specialisation, or modularity itself -- we measure accuracy,
tolerance to module lesion and sensor failure, and inference cost.
\todo{one-sentence headline result, written after the numbers exist}
\end{abstract}

\begin{IEEEkeywords}
modular neural networks, brain-inspired architectures, multimodal fusion,
embodied navigation, conditional computation, damage tolerance
\end{IEEEkeywords}

% ---------------------------------------------------------------------------
\section{Introduction}

A single feedforward network applied to concatenated inputs is the default
design for multimodal systems. It is also the design the brain conspicuously
does not use. Cortex is regionally specialised: primary visual cortex is
retinotopic and locally connected; posterior parietal cortex computes
relations between locations; temporal cortex handles sequence and meaning;
prefrontal cortex integrates and decides. Sensory input does not enter through
one door but through modality-specific thalamic nuclei -- the lateral
geniculate for vision, the medial geniculate for audition, the ventral
posterior nuclei for somatosensation.

Whether this organisation confers a computational advantage, as opposed to
being a developmental or metabolic constraint, is an empirical question that
architectural work in machine learning rarely tests directly. Systems
described as ``brain-inspired'' frequently reproduce the \emph{names} of
cortical regions while implementing identical subnetworks, in which case the
biological framing is decorative and any result is attributable to capacity
rather than structure.

This paper takes the opposite approach. We make the modules structurally
different, hold parameters constant, and remove one property at a time.

\textbf{Contributions.}
\begin{enumerate}
  \item A modular architecture in which each cortical module carries a
        \emph{distinct} inductive bias matched to the structure of the signal
        routed to it, rather than a shared block under different names
        (Section~\ref{sec:arch}).
  \item A modality-keyed thalamic relay including a somatosensory pathway to
        parietal cortex, which lets a proprioceptive sensor participate in
        spatial reasoning without passing through vision
        (Section~\ref{sec:thalamus}).
  \item \textsc{GridNav}, an embodied navigation benchmark constructed so that
        no single modality suffices, with an exact shortest-path oracle and a
        published per-modality ceiling (Section~\ref{sec:gridnav}).
  \item An ablation protocol that isolates routing, specialisation and
        modularity separately, under parameter matching and paired seeds
        (Section~\ref{sec:setup}).
\end{enumerate}

% ---------------------------------------------------------------------------
\section{Related Work}

\todo{write this section. It must cover, and honestly position against:}
\begin{itemize}
  \item \emph{Mixture-of-experts and conditional computation.} Our routing is
        fixed and modality-keyed rather than learned per-token; the closest
        comparison is to hard-routed MoE. State plainly that we do not claim
        novelty for conditional computation itself.
  \item \emph{Modular and compositional networks} -- neural module networks,
        recurrent independent mechanisms, capsule-style routing.
  \item \emph{Multimodal fusion} -- early/late/intermediate fusion, and where
        frontal convergence sits among them.
  \item \emph{Brain-inspired architectures}, including work that uses cortical
        naming without structural differentiation. This is the gap the paper
        claims.
  \item \emph{Catastrophic forgetting and modularity}, as the motivation for
        why separable modules might help.
\end{itemize}

% ---------------------------------------------------------------------------
\section{Architecture}
\label{sec:arch}

\subsection{Common representational width}

Every module accepts and emits a vector of fixed width $D$ ($D=256$
throughout). Fixing $D$ makes modules interchangeable at the wiring level,
which is what allows an ablation to replace one module with another and change
nothing else. All derived widths follow from $D$.

Raw inputs reach this common space through per-modality projections. Vision,
language and audio use frozen pretrained encoders (ResNet-18, MiniLM,
Whisper-tiny) followed by a trainable linear projection; the backbones are
never updated, so no result can be attributed to one variant obtaining better
features. The proprioceptive modality has no pretrained backbone and does not
want one -- a rangefinder reading is already a measurement -- so its encoder is
a small trainable MLP. This matters for interpretation: navigation performance
on that modality is not inherited from a large pretrained model.

\subsection{Modality-keyed thalamic relay}
\label{sec:thalamus}

Input is not funnelled through a single entry point. Each modality is relayed
by its own nucleus to a specific cortical target:
\begin{center}
\begin{tabular}{lll}
\toprule
Modality & Nucleus & Target \\
\midrule
image  & LGN      & occipital \\
audio  & MGN      & temporal  \\
text   & pulvinar & temporal  \\
sensor & VPL      & parietal  \\
\bottomrule
\end{tabular}
\end{center}

Each nucleus applies a learned multiplicative gate,
$\mathrm{relay}(x)\odot\sigma(\mathrm{gate}(x))$, rather than forwarding
unchanged; thalamic relay cells gain-control what passes rather than acting as
wires.

Two properties of this table are load-bearing. First, vision reaches parietal
cortex only \emph{via} occipital, mirroring the dorsal stream; a direct
LGN--parietal path is rejected by construction. Second, somatosensory input
does reach parietal directly, which is anatomically correct and is what allows
an embodied agent's body sense to inform spatial reasoning without being
rendered into an image first.

The text pathway is a deliberate simplification and we flag it rather than
bury it: reading in humans enters through vision, and no thalamic nucleus
relays orthography to temporal cortex. We route text directly because the
alternative -- rendering text to pixels -- would confound the language pathway
with the visual one.

\subsection{Functionally distinct cortical modules}

This is the component that distinguishes the work from architectures that use
cortical vocabulary without cortical structure. The four modules do not share
an implementation:

\begin{itemize}
  \item \textbf{Occipital} uses \emph{grouped} connectivity: the representation
        is partitioned into bands and each hidden unit reads one band, giving
        local receptive fields, followed by a mixing layer that integrates
        across bands. A dense layer discards this locality.
  \item \textbf{Parietal} reads the representation as a set of positions and
        applies self-attention across them. A spatial relation is a comparison
        \emph{between} locations, an operation an elementwise MLP cannot
        express regardless of width.
  \item \textbf{Temporal} reads the representation as an ordered sequence and
        applies a gated recurrent unit. Both modalities routed here, language
        and audio, are sequential; an MLP is order-invariant and discards that.
  \item \textbf{Frontal} is the sole convergence point. It receives a variable
        number of $D$-wide streams and fuses them before deciding.
\end{itemize}

Because these operators differ, ablating a module removes a \emph{capability}
rather than merely a block of parameters -- which is what makes the ablation in
Section~\ref{sec:results} interpretable.

\subsection{Convergence and fusion}

Frontal cortex receives from occipital directly \emph{and} via parietal, so on
the visual pathway it receives two streams simultaneously. With a fixed width
$D$ this is a shape conflict, resolved by fusing before the module body:
\[
n \times D \;\longrightarrow\; \mathrm{fuse} \;\longrightarrow\; D
   \;\longrightarrow\; \mathrm{module} \;\longrightarrow\; D .
\]
Streams occupy \emph{fixed positional slots}. A modality that is absent, or a
module that is lesioned, contributes zeros to its own slot rather than shifting
its neighbours. Without fixed slots the fusion weights would mean different
things depending on which modalities happened to be present, and lesion
experiments would be uninterpretable. We use concatenation rather than
summation so that provenance is preserved; summation is available and is
compared in the ablations.

\subsection{Basal-ganglia gating}

Frontal cortex's decision passes a go/no-go gate before reaching the output.
A shared striatal representation feeds a direct (Go) and an indirect (NoGo)
pathway, and the gate is $\sigma(\mathrm{go}-\mathrm{nogo})$, applied
multiplicatively.

We note one non-obvious implementation constraint, because it silently
disabled the module in our first version. Initialising both pathway weight
matrices to zero leaves the gradient reaching the striatum proportional to
their difference, which stays near zero because the two pathways receive
exactly opposite gradients; the striatal parameters then never train. Small
random initialisation breaks the symmetry while biases keep the gate open at
the start.

% ---------------------------------------------------------------------------
\section{The GridNav Benchmark}
\label{sec:gridnav}

Multimodal benchmarks frequently permit a single modality to carry most of the
signal, in which case fusion is not tested. \textsc{GridNav} is constructed to
prevent that.

An agent occupies a randomly walled $12\times12$ occupancy grid with one goal.
Supervision is exact: a breadth-first distance field from the goal defines the
optimal move, so there is no reward shaping and no exploration variance. Where
several moves are equally optimal -- which occurs in \todo{X}\% of states --
all are accepted, and accuracy is measured as whether the chosen action is in
the optimal set. Scoring against one arbitrarily chosen label would penalise
correct behaviour.

Each modality is deliberately partial:
\begin{itemize}
  \item \textbf{Vision} is an agent-centric $5\times5$ window. The goal falls
        inside it in only \todo{X}\% of states, so vision reports obstacles but
        rarely direction.
  \item \textbf{Proprioception} is eight rangefinder beams giving normalised
        distance to the nearest wall and a contact flag. It carries no goal
        information at all.
  \item \textbf{Language} is a coarse instruction (``the goal lies north-east,
        far away''). It carries no obstacle information at all.
\end{itemize}

Direction without obstacles walks into walls; obstacles without direction have
nowhere to go. We report the single-modality ceilings in
Table~\ref{tab:sufficiency} as evidence that fusion is required rather than
asserting it.

The monolithic baseline receives the concatenation of the same three
observations at the same resolution, so it is disadvantaged in structure only,
never in information.

% ---------------------------------------------------------------------------
\section{Experimental Setup}
\label{sec:setup}

\subsection{Tasks}
\textsc{GridNav} (above) exercises every module. We add two single-modality
controls: CIFAR-100 (vision) and 20 Newsgroups (language). We include these
\emph{expecting no modular advantage} -- with one modality there is nothing to
route -- because a paper that reported only the setting favourable to its
hypothesis would not be evidence.

\subsection{Variants}
Each baseline removes exactly one property:
\begin{itemize}
  \item \textbf{--routing}: every modality is broadcast to every module.
  \item \textbf{--specialisation}: four separate modules, all the plain block.
  \item \textbf{--modularity}: one module instance reused for every stream.
  \item \textbf{Monolithic}: a single MLP over the concatenated input.
\end{itemize}
Comparing only against the monolithic model would confound three changes at
once and could not identify which property mattered.

\subsection{Protocol}
All variants are matched to the modular model's trainable parameter count
(within \todo{X}\%; achieved counts in Table~\ref{tab:cost}) and share the
optimiser, schedule, batch size and epoch budget. Frozen encoder features are
computed once and cached, so every variant consumes identical inputs.

Model selection uses a validation split held out from training; the test split
is read exactly once per run. Each configuration runs over $n$ seeds, and
seeds are \emph{paired} across variants -- seed $k$ fixes the same split and
initialisation draw for every model -- so comparisons use a paired $t$-test and
report Cohen's $d_z$ alongside $p$.

\subsection{Failure modes}
We measure two distinct failures. \emph{Module lesion} disables a cortical
module, which is the biological analogue. \emph{Sensor failure} removes an
input modality, which is the engineering case for an embodied agent whose
camera is occluded or whose rangefinder fails. These are different experiments
and we keep them separate.

% ---------------------------------------------------------------------------
%% IMPORTANT: the tables below are generated. Do not type numbers into this
%% section. Write the interpreting sentences ONLY after reading results.json,
%% and if the modular model loses on a task, say so -- the ablations are
%% informative either way, and a negative multimodal result is publishable.
% ---------------------------------------------------------------------------
\section{Results}
\label{sec:results}

\input{generated/table_main}

\todo{Interpret Table~\ref{tab:main}. Expect no modular advantage on the
single-modality controls; state that plainly.}

\input{generated/table_significance}

\todo{Interpret Table~\ref{tab:significance}. Which single removed property
costs the most? That is the paper's actual finding.}

\subsection{Damage tolerance}

\input{generated/table_damage}

\begin{figure}[t]
  \centering
  \includegraphics[width=\columnwidth]{generated/damage.pdf}
  \caption{Accuracy retained on \textsc{GridNav} under module lesion and
           sensor failure.}
  \label{fig:damage}
\end{figure}

\todo{Interpret. The claim to test is graceful rather than catastrophic
degradation relative to the monolithic model.}

\subsection{Inference cost}

\input{generated/table_cost}

\todo{Interpret. Routing should skip modules no modality addresses; report the
saving, or report that there is none.}

% ---------------------------------------------------------------------------
\section{Discussion}

\todo{Written after results.}

\subsection{Limitations}

We state these regardless of outcome.

\begin{itemize}
  \item \textbf{Scale.} Modules are on the order of $10^6$ parameters and tasks
        are small enough to train on CPU. Whether the effect survives at scale
        is untested.
  \item \textbf{Routing is fixed, not learned.} The modality--module map is
        specified by anatomy, not discovered. A learned router is the obvious
        comparison and is not run here.
  \item \textbf{\textsc{GridNav} is synthetic.} The oracle is exact and the
        environment is fully observable to the generator, which is what makes
        supervision clean but also what makes it easier than a real robot's
        world.
  \item \textbf{Borrowed perception.} Three of four modalities depend on frozen
        pretrained encoders. The architecture is evaluated on what it does with
        features, not on learning them.
  \item \textbf{No temporal credit assignment.} \textsc{GridNav} is single-step
        imitation of a shortest-path oracle, not sequential control. Nothing
        here demonstrates closed-loop navigation.
  \item \textbf{Audio is unexercised.} The pathway exists and is verified, but
        no task in this evaluation routes audio.
\end{itemize}

% ---------------------------------------------------------------------------
\section{Conclusion}

\todo{Written after results.}

% ---------------------------------------------------------------------------
\bibliographystyle{IEEEtran}
\bibliography{refs}

\end{document}
